\chapter{The theorem and the argument}
\label{chap:introduction}

A family of thin tubes can have small union because many of its members
pass through the same points. The sticky Kakeya problem asks whether this
overlap can persist at every scale when the family remains suitably spread
out after rescaling. In three dimensions, Wang and Zahl showed that it
cannot. This blueprint follows the finite-scale argument and its
formalization in Lean.

Our endpoint is Theorem~5.2 of \cite{WZ2}: a nonempty family of
$\delta$-tubes satisfying the Convex Wolff Axioms at every scale, with a
dense shading, has union of volume at least $\delta^\eps$. The allowed
nonconcentration and density losses depend on $\eps$ and tend to zero
with it. The hypotheses and their quantifier order are given explicitly
in Theorem~\ref{thm:sticky}. This is the sticky estimate used in the
project; the subsequent passage to the general three-dimensional Kakeya
conjecture lies beyond the present account.

We write WZ1 for \emph{Sticky Kakeya sets and the sticky Kakeya
conjecture}~\cite{WZ1}, and WZ2 for \emph{The Assouad dimension of Kakeya
sets in $\R^3$}~\cite{WZ2}. The numbering attached to statements below
is the numbering of the cited paper, in addition to this document's own
numbering. We use the October 2025 arXiv version of WZ1 and the April
2025 arXiv version of WZ2. Free author manuscripts are used for the other
references. A change of numbering between manuscripts and published
articles is recorded where it matters.

\section{The continuum statement}

\begin{definition}[Sticky sets; WZ1, Definitions 1.1--1.2]
\label{def:continuum-sticky}
A compact set $K\subset\R^n$ is a sticky Kakeya set if there is a
family $L$ of affine lines of packing dimension $n-1$, containing a
line in every direction, such that $K\cap\ell$ contains a unit
segment for every $\ell\in L$. More generally, $K$ is
$\eta$-close to being sticky if such a family can be chosen with
packing dimension at most $n-1+\eta$. Packing dimension is measured
in the usual metric on the space of affine lines, equivalently in
bounded line-parameter charts.
\end{definition}

\begin{theorem}[The sticky Kakeya theorem; WZ1, Theorem 1.1]
\label{thm:wz1-continuum}
\uses{def:continuum-sticky}
Every sticky Kakeya set in $\R^3$ has Hausdorff dimension $3$.
More quantitatively, for every $\eps>0$ there is $\eta>0$ such that
every Kakeya set in $\R^3$ that is $\eta$-close to being sticky has
Hausdorff dimension at least $3-\eps$; this is WZ1 Theorem 1.1$'$.
\end{theorem}

This is the continuum statement motivating the finite problem.
Proposition~\ref{prop:wz1-dimension-reduction} explains the role of
the WZ1 critical exponent. The final Lean declaration linked in this
blueprint is the all-scale-CWA finite estimate, not a separate
formalization of the continuum formulation above.

\section{Why grain structure leads to a contradiction}

Suppose the sticky estimate fails. One can then choose an exponent
$\sigma>0$ and families whose shaded union has volume approximately
$\delta^\sigma$, and for which this exponent cannot be improved. Wolff's
hairbrush estimate places $\sigma$ strictly below $1$. Extremality has
two consequences. First, both a coarse family and the rescaled fine
family inside a typical coarse tube must themselves be extremal.
Second, their multiplicities must be close to their largest permissible
values. These observations make it possible to carry information from
one scale to another.

The multilinear Kakeya estimate forces the tubes through a typical point
to lie approximately in a plane. Comparing these planes at successive
scales produces a Lipschitz plane map and local grain structure. After
a suitable choice of coordinates, the horizontal slices of the shaded
union also have global grains, whose direction is described by a
Lipschitz function $f(z)$. A projection theorem, proved using radial
projections, improves this function to $C^2$.

The next step is to show that $f$ cannot be nearly constant for too long.
Otherwise a thin rectangular prism would contain too much of the tube
family, contrary to the Convex Wolff Axioms. After an anisotropic
rescaling, one obtains
\[
  1\leq |f'(z)|\leq 2,\qquad |f''(z)|\leq 1/100.
\]
Global grains make the twisted projection
\[
  \pi_f(x,y,z)=(x+f(z)y,z)
\]
small. On the other hand, the image of a line under $\pi_f$ is a curve
in a cinematic family. A maximal estimate for these curves, combined
with an iteration over scales, makes the same projection large. The
two estimates contradict each other when the losses are chosen in the
correct order.

\section{Reading the formalization}

The Lean links attached to statements lead to their declarations in the
repository. Definitions, theorem statements, and proved implications
play different roles. When a paper result is used only in a normalized
or restricted form, the text states that form and explains the
difference in a remark. In particular, a link to a proposition-valued
definition does not by itself assert that the proposition has been
proved. Completion marks are reserved for the mathematical formulation
described at that node.

The paper argument graph records the mathematical uses described in the
proofs. A separate \emph{Lean proof dependencies} view, linked from the
project homepage, is extracted from the compiled sticky theorem. It
lists the actual Lean versions reached from that theorem, together with
source paths witnessing the dependencies. In particular, a paper node
can link several proved versions of which only one is used, or an
interface that is not itself a proof. The two graphs answer different
questions; the paper graph's colours are the manual completion marks.

The paper graph omits routine implementation lemmas. Most of the numbered
results in WZ1's structural and projection arguments and WZ2's sticky
argument are included, together with the supporting numbered results
from the projection and cinematic estimates. The correspondence chapter
records the scope and the adjustments that affect the meaning of a
statement.
